\documentclass[10pt,a4paper]{amsart}
\usepackage[margin=19mm]{geometry}
\usepackage{amsmath,amssymb,mathtools}
\usepackage[scr=rsfs,cal=euler]{mathalfa}
\usepackage{xfrac}
\usepackage[unicode,hidelinks,bookmarks=false]{hyperref}
\newcommand{\ie}{i.e.\ }
\newcommand{\cf}{cf.\ }
\newcommand{\resp}{respectively\ }
\newcommand{\Subsection}{\subsection}
\providecommand{\nobreakdash}{\nobreak\hbox{-}}
\setlength{\emergencystretch}{2em}
\multlinegap0pt
\usepackage{amssymb}
\usepackage{bm}
\usepackage[shortlabels]{enumitem}
\usepackage{xcolor}
\usepackage{tikz}
\newtheorem{proposition}{Proposition}
\newcommand{\R}{{\mathbb R}}
\title{\textbf{A Moser-Type Construction for the Liouville Equation}}
\author{A. Borz\`i \and M. Caponigro \and A. Vicari}
\date{Source: arXiv:2603.21850v2 (CC BY 4.0)}
\begin{document}
\maketitle

\noindent\emph{Source and licence.} \textbf{A Moser-Type Construction for the Liouville Equation}. Version arXiv:2603.21850v2 (CC BY 4.0), distributed by arXiv under the Creative Commons Attribution 4.0 International licence, http://creativecommons.org/licenses/by/4.0/, as recorded in the arXiv licence metadata for this exact version and shown on the abstract page https://arxiv.org/abs/2603.21850v2. Full licence notice: \texttt{licenses/arxiv-2603.21850-CC-BY-4.0.txt}.\par\smallskip

\noindent\textit{Editorial scope.} Counted unit: the article's proposition prop:equimeasurable (source0.tex lines 305--347). The article's complete original proof of it is delivered with it (source0.tex lines 349--404, one \texttt{proof} environment of 56 lines). No mathematics is written, repaired or re-derived: the statement and the proof are the article's own lines, in the article's own notation, and no retained line is substituted or re-flowed.  Retained uncounted as the source-local context the proof needs: lines 153--156; lines 291--296.  Nothing was omitted from any retained span, so the package carries no omission record.  No retained line contains a citation, so the wrapper carries no bibliography and no bibliography entry.  No retained line contains a figure environment, an image-inclusion command or drawing code, so no image is delivered and the package declares no asset dependency.  Editorial formatting only, in addition: an AMS \texttt{amsart} preamble that restates the article's own declarations and macros used by the retained lines, namely the environments \texttt{proposition} on separate counters, together with the standard packages those lines need.  The retained environments print in this document as Proposition 1 in order of appearance, whereas the article prints the same items with the numbering of its own full document.  The complete native source of the article as distributed in the arXiv e-print for this exact version is bundled unchanged as \texttt{sources/196978/source0.tex}.
\begin{equation}\label{eqLiouville}
\partial_t \rho(x,v,t) + v\cdot\nabla_x \rho(x,v,t)  
+  \nabla_v \cdot \big( a(x,v,t) \, \rho(x,v,t)  \big)= 0, 
\end{equation}

\begin{equation}\label{eq:hamtimedep}
\begin{cases}
\dot x &=v,\\
\dot v &=-\nabla_x V(x,t),
\end{cases}
\end{equation}

\begin{proposition}
\label{prop:equimeasurable}
Let $\rho(x,t)$ be the solution 
of the Liouville equation
\[
\partial_t\rho + v\cdot\nabla_x\rho - \nabla_x V(x,t)\cdot\nabla_v\rho = 0,
\qquad
\rho(\cdot,0)=f,
\]
such that 
\[
\rho(\cdot,1)=g .
\]
Then $f$ and $g$ are equimeasurable, namely
\[
        \int_{\Omega} \Phi(f(x,v))\,dx\,dv
        =
        \int_{\Omega} \Phi(g(x,v))\,dx\,dv
\]
for any measurable function $\Phi:\R \to \R$.
%Let $\Omega=X\times V$ and consider the Hamiltonian system
%\[
%\dot x=v, \qquad \dot v=-\nabla_x V(x,t),
%\]
%where $V$ is sufficiently smooth. Let $\Psi_t:\Omega\to\Omega$ denote the
%corresponding Hamiltonian flow map, and let $\rho$ solve the Liouville equation
%\[
%\partial_t\rho + v\cdot\nabla_x\rho - \nabla_x V(x,t)\cdot\nabla_v\rho = 0,
%\qquad
%\rho(\cdot,0)=f .
%\]
%Assume that
%\[
%\rho(\cdot,1)=g .
%\]
%Then $f$ and $g$ are equimeasurable, namely
%\[
%|\{(x,v)\in\Omega:f(x,v)>\lambda\}|
%=
%|\{(x,v)\in\Omega:g(x,v)>\lambda\}|
%\]
%for every $\lambda\in\R$.
\end{proposition}

\begin{proof}
The phase-space vector field of~\eqref{eq:hamtimedep} is divergence-free 
since $\nabla_x\cdot v+\nabla_v\cdot(-\nabla_x V)=0$.
As a consequence, the associated flow is volume-preserving in phase space.
If \(\Psi_t\) denotes the Hamiltonian flow, then the solution of~\eqref{eqLiouville} is
\[
        \rho(t)=\rho(0)\circ \Psi_t^{-1}.
\]
The Hamiltonian flow $\Psi_t$ preserves phase-space volume
(Liouville theorem), that is,
\[
\det D\Psi_t(x,v)=1 .
\]
Along characteristics, the density satisfies
\[
\frac{d}{dt}\rho(\Psi_t(x,v),t)=0,
\]
hence $\rho(\Psi_t(x,v),t)=f(x,v)$. At time $t=1$ we obtain $g(x,v)=f(\Psi_1^{-1}(x,v))$. 

Fix $\lambda\in\R$ and define the level sets
\[
E_f(\lambda)
=
\{z\in\Omega:f(z)>\lambda\},
\qquad
E_g(\lambda)
=
\{z\in\Omega:g(z)>\lambda\}.
\]
Since $g(z)=f(\Psi_1^{-1}(z))$, we have
\[
z\in E_g(\lambda)
\iff
\Psi_1^{-1}(z)\in E_f(\lambda),
\]
that is,
\[
E_g(\lambda)=\Psi_1(E_f(\lambda)).
\]

Because $\Psi_1$ is volume-preserving,
\[
|E_g(\lambda)|
=
|\Psi_1(E_f(\lambda))|
=
|E_f(\lambda)|.
\]
Therefore
\[
|\{g>\lambda\}|
=
|\{f>\lambda\}|,
\]
and the conclusion follows.
\end{proof}

\end{document}
